%============================================================
% CH1 · P4 · PROBLEM 2 — BUILD ARRAY FROM PERMUTATION (LeetCode 1920)
%============================================================
\subsection{Problem 1.4.2: Build Array from Permutation}
\problemlabel{1.4.2}{ssec:ch1-build-permutation}

\problemcard{1.4.2}{Build Array from Permutation}{LeetCode 1920}{https://leetcode.com/problems/build-array-from-permutation/}{\exeasy\ \lc{Easy}}{Index arithmetic $i \to a[a[i]]$; in place by storing old $+\, n \cdot$new in one cell}

\subsubsection*{1. Problem statement}

You are given a \textbf{zero-based permutation} \texttt{nums}: an array of the $n$ distinct integers $0, 1, \dots, n-1$ in some order. Build the array \texttt{ans} of the same length with
\[
\texttt{ans}[i] = \texttt{nums}[\,\texttt{nums}[i]\,] \qquad \text{for every } 0 \le i < n,
\]
and return it. \textbf{Follow-up:} can you do it without extra space, i.e.\ with $\bigO{1}$ memory besides the input array?

\textbf{Constraints:} $1 \le n \le 1000$, \; $0 \le \texttt{nums[i]} < n$, \; all values distinct.

\begin{table}[H]
\centering\small
\begin{tabular}{@{}lll@{}}
\toprule
\textbf{nums} & \textbf{Output} & \textbf{For example} \\
\midrule
$[0,2,1,5,3,4]$ & $[0,1,2,4,5,3]$ & $\texttt{ans}[3] = \texttt{nums}[\texttt{nums}[3]] = \texttt{nums}[5] = 4$ \\
$[5,0,1,2,3,4]$ & $[4,5,0,1,2,3]$ & $\texttt{ans}[0] = \texttt{nums}[5] = 4$ \\
\bottomrule
\end{tabular}
\end{table}

\begin{diagrambox}
\begin{tikzpicture}[x=1cm, y=1cm]
    \dpindex{ix}{-1.3}{index}{0}{5}
    \dparray{a}{1.5}{nums}{0/acell, 2/acell, 1/acell, 5/aread, 3/acell, 4/aread}
    \dparray{b}{-0.6}{ans}{0/akeep, 1/akeep, 2/akeep, 4/awrite, 5/akeep, 3/akeep}
    \draw[dep] (a3.north) to[bend left=35] node[lbl, above]{$\texttt{nums}[3] = 5$: go to cell 5} (a5.north);
    \draw[dep] (a5.south) to[bend left=10] node[lbl, right, pos=0.6]{read 4} (b3.north);
    \node[note, anchor=west, text width=3.4cm, align=left] at (5.8,0.45) {two jumps: index $\to$ value $\to$ value at that index};
\end{tikzpicture}
\figcap{Sample 1. The cell value is used as an index: $3 \to \texttt{nums}[3] = 5 \to \texttt{nums}[5] = 4$.}
\end{diagrambox}

\subsubsection*{2. The problem model}
\begin{modelbox}[Problem Model --- Build Array from Permutation]
\renewcommand{\arraystretch}{1.35}
\begin{tabularx}{\linewidth}{@{}>{\bfseries}l X@{}}
Input & a permutation $p$ of $\{0, \dots, n-1\}$, stored as $p[i] = \texttt{nums}[i]$ \\
Computed & the array $p \circ p$: position $i$ receives $p(p(i))$ (``apply the map twice'') \\
Objects & the $n$ positions; values double as indices, so every value is a valid index \\
Size & $n \le 1000$: $\bigO{n}$ is trivial; the real constraint is the follow-up --- $\bigO{1}$ extra memory, so the old values must survive while the array is being overwritten \\
\end{tabularx}
\end{modelbox}

\begin{diagrambox}
\begin{tikzpicture}[x=1cm, y=1cm]
    \foreach \i/\x/\y in {0/0/0, 1/1.8/0.7, 2/1.8/-0.7, 3/4.2/0.8, 5/5.6/0, 4/4.2/-0.8}{
        \node[gv] (v\i) at (\x,\y) {\i};
    }
    \draw[gde] (v0) to[loop left, looseness=6] (v0);
    \draw[gde] (v1) to[bend left=25] (v2); \draw[gde] (v2) to[bend left=25] (v1);
    \draw[gde] (v3) -- (v5); \draw[gde] (v5) -- (v4); \draw[gde] (v4) -- (v3);
    \node[note, anchor=west, text width=5.4cm, align=left] at (6.6,0)
        {arrow $i \to \texttt{nums}[i]$. A permutation is a set of disjoint cycles: $(0)$, $(1\,2)$, $(3\,5\,4)$.\\[3pt]
        $\texttt{ans}[i]$ is the vertex \emph{two} arrows ahead of $i$: $\texttt{ans}[3] = 4$, $\texttt{ans}[1] = 1$.};
\end{tikzpicture}
\figcap{The permutation of Sample 1 as arrows. Building \texttt{ans} = moving two steps along the cycles.}
\end{diagrambox}

\subsubsection*{3. Key observations}

\begin{itemize}[leftmargin=*]
    \item \textbf{With a second array the problem is one line:} \texttt{ans[i] = nums[nums[i]]}. The difficulty is entirely in the follow-up.
    \item \textbf{The tempting in-place version is wrong.} \texttt{nums[i] = nums[nums[i]]} for $i = 0, 1, \dots$ reads cells that were already overwritten. On Sample 1 it produces $[0,1,1,4,4,4]$ instead of $[0,1,2,4,5,3]$: at $i = 2$ it reads $\texttt{nums}[1]$, which already holds the \emph{new} value 1 instead of the old 2. No iteration order fixes this in general, because the dependencies form cycles.
    \item \textbf{Every value is smaller than $n$.} So a cell can carry two numbers in $[0, n)$ at once: $\text{old} + n \cdot \text{new}$ (Tool~4 of the theory). The old value is always $\texttt{nums}[j] \bmod n$, whether cell $j$ has been rewritten yet or not.
    \item \textbf{No overflow:} the largest encoded value is $(n-1) + n(n-1) = n^2 - 1 < 10^6$, comfortably inside \texttt{int}.
\end{itemize}

\begin{diagrambox}
\begin{tikzpicture}[x=1cm, y=1cm]
    \dparray{w}{0.9}{naive}{0/acell, 1/awrite, 1/aout, 4/awrite, 4/aout, 4/aout}
    \dparray{r}{-0.3}{correct}{0/akeep, 1/akeep, 2/akeep, 4/akeep, 5/akeep, 3/akeep}
    \draw[dep, warningred] ([yshift=1mm]w1.north) to[bend left=50] node[lbl, above, text=warningred]{read at $i=2$: new value!} ([yshift=1mm]w2.north);
    \node[note, anchor=west, text width=3.7cm, align=left] at (5.8,0.3) {grey cells are wrong: they were computed from cells that had already changed};
\end{tikzpicture}
\figcap{Counter-example to \texttt{nums[i] = nums[nums[i]]} in place on Sample 1.}
\end{diagrambox}

\subsubsection*{4. Approaches}

\begin{bruteforcebox}
\textbf{Idea.} Write the answer into a second array, so the old values stay readable.
\begin{lstlisting}[style=cpp]
#include <bits/stdc++.h>
using namespace std;

class Solution {
public:
    vector<int> buildArray(vector<int>& nums) {
        int n = nums.size();
        vector<int> ans(n);                    // a second array: old values stay readable
        for (int i = 0; i < n; i++) ans[i] = nums[nums[i]];
        return ans;
    }
};

int main() {                                   // input: n, then a permutation of 0..n-1
    int n;
    if (!(cin >> n)) return 0;
    vector<int> a(n);
    for (auto& x : a) cin >> x;
    vector<int> r = Solution().buildArray(a);
    for (int i = 0; i < n; i++) cout << r[i] << (i + 1 < n ? ' ' : '\n');
}
\end{lstlisting}
\begin{complexitybox}
\textbf{Time:} $\bigO{n}$. \textbf{Space:} $\bigO{n}$ extra. \ac{} for the base problem; it fails only the follow-up, and it is the reference for the stress test.
\end{complexitybox}
\end{bruteforcebox}

\begin{algorithmbox}
\textbf{Idea.} Pass 1: for every $i$, read the old value $j = \texttt{nums}[i] \bmod n$, read the old value of cell $j$ as $\texttt{nums}[j] \bmod n$ (that is the new value of cell $i$), and add $n$ times it to cell $i$. Pass 2: divide every cell by $n$ --- the old ``low digit'' disappears, the new ``high digit'' remains.
\begin{lstlisting}[style=cpp]
#include <bits/stdc++.h>
using namespace std;

class Solution {
public:
    vector<int> buildArray(vector<int>& nums) {
        int n = nums.size();
        // pass 1: cell i keeps old = nums[i] % n, adds n * new (new = old value of nums[old])
        for (int i = 0; i < n; i++)
            nums[i] += n * (nums[nums[i] % n] % n);   // % n reads the OLD value of any cell
        // pass 2: drop the old value, keep the new one
        for (int i = 0; i < n; i++)
            nums[i] /= n;
        return nums;
    }
};

int main() {                                          // input: n, then a permutation of 0..n-1
    int n;
    if (!(cin >> n)) return 0;
    vector<int> a(n);
    for (auto& x : a) cin >> x;
    vector<int> r = Solution().buildArray(a);
    for (int i = 0; i < n; i++) cout << r[i] << (i + 1 < n ? ' ' : '\n');
}
\end{lstlisting}
\begin{complexitybox}
\textbf{Time:} $\bigO{n}$ --- two passes. \textbf{Space:} $\bigO{1}$ extra: the answer is returned in the input array. \textbf{Verdict:} \ac{}, and it meets the follow-up.
\end{complexitybox}
\end{algorithmbox}

\begin{invariantbox}[Why the encoding is correct]
\textbf{Invariant of pass 1.} Before iteration $i$, every cell $j < i$ holds $\text{old}_j + n \cdot \text{new}_j$ and every cell $j \ge i$ holds $\text{old}_j$, where $\text{new}_j = \text{old}_{\text{old}_j}$. In both cases $\texttt{nums}[j] \bmod n = \text{old}_j$, because $0 \le \text{old}_j < n$ and adding a multiple of $n$ does not change the remainder.

\textbf{Step.} Iteration $i$ computes $\texttt{nums}[i] \bmod n = \text{old}_i$ (cell $i$ is still unencoded) and then $\texttt{nums}[\text{old}_i] \bmod n = \text{old}_{\text{old}_i} = \text{new}_i$ by the remark above, and adds $n \cdot \text{new}_i$ to cell $i$. The invariant holds for $i+1$.

\textbf{Pass 2.} Since $0 \le \text{old}_i < n$, $\lfloor (\text{old}_i + n\,\text{new}_i)/n \rfloor = \text{new}_i = \texttt{nums}[\texttt{nums}[i]]$ (old values). \hfill$\blacksquare$
\end{invariantbox}

\begin{alternativebox}
\textbf{Shift the values along each cycle (no arithmetic encoding).} On a cycle $c_0 \to c_1 \to \dots \to c_{L-1} \to c_0$ (where $\texttt{nums}[c_t] = c_{t+1}$), the new value of $c_t$ is $c_{t+2}$, which is the \emph{old value of the next cell} $c_{t+1}$. So each cycle just shifts its values one step back --- exactly the cycle-walking idea of Pattern~3 (Tool~3). Walk the cycle from its first cell $s$; each cell takes the old value of the next one, and the last cell takes the saved old value of $s$. Mark rewritten cells by storing $\sim\!v = -v - 1$ (negative, even for $v = 0$), and undo the marks at the end.
\begin{lstlisting}[style=cpp]
#include <bits/stdc++.h>
using namespace std;

class Solution {
public:
    vector<int> buildArray(vector<int>& nums) {
        int n = nums.size();
        for (int s = 0; s < n; s++) {
            if (nums[s] < 0) continue;            // already rewritten (stored as ~value)
            // cycle s -> p(s) -> p(p(s)) -> ... : the new value of a cell is the old value of
            // the next cell of the cycle, so shift the values one step along the cycle
            int first = nums[s], i = s;
            while (nums[i] != s) {                // stop at the cell that points back to s
                int j = nums[i];                  // next cell of the cycle
                nums[i] = ~nums[j];               // new value of i = old value of j, marked
                i = j;
            }
            nums[i] = ~first;                     // the last cell takes the old value of s
        }
        for (int i = 0; i < n; i++) nums[i] = ~nums[i];   // remove the marks (~~x = x)
        return nums;
    }
};

int main() {                                      // input: n, then a permutation of 0..n-1
    int n;
    if (!(cin >> n)) return 0;
    vector<int> a(n);
    for (auto& x : a) cin >> x;
    vector<int> r = Solution().buildArray(a);
    for (int i = 0; i < n; i++) cout << r[i] << (i + 1 < n ? ' ' : '\n');
}
\end{lstlisting}
Every cell is written once in the cycle walk and once when the marks are removed: $\bigO{n}$ time, $\bigO{1}$ memory, and no assumption that $n^2$ fits in the type --- it only needs one free bit (the sign) per cell. Why $\sim$ and not the minus sign? $-0 = 0$, so the cell holding 0 could not be marked.

\begin{diagrambox}
\begin{tikzpicture}[x=1cm, y=1cm]
    \dparray{c}{0}{old}{0/acell, 2/acell, 1/acell, 5/aread, 3/aread, 4/aread}
    \draw[dep] (c5.north) to[bend right=40] node[lbl, above]{to cell 3} (c3.north);
    \draw[dep] (c4.south) to[bend left=50] (c5.south);
    \draw[dep, dashed] (c3.south) to[bend right=50] (c4.south);
    \node[note, anchor=west, text width=3.6cm, align=left] at (5.9,0) {cycle $3 \to 5 \to 4$: cell 3 gets old $a_5 = 4$, cell 5 gets old $a_4 = 3$, cell 4 gets the saved old $a_3 = 5$ (dashed)};
\end{tikzpicture}
\figcap{The cycle $(3\,5\,4)$ of Sample 1: its values $5, 4, 3$ (in cycle order) shift by one position.}
\end{diagrambox}
\end{alternativebox}

\begin{tricksbox}
\begin{itemize}[leftmargin=*]
    \item \textbf{Bit packing instead of base $n$:} values below $1024$ fit in 10 bits, so \texttt{nums[i] |= (nums[nums[i] \& 1023] \& 1023) << 10}, then \texttt{nums[i] >>= 10}. Same idea, shifts instead of multiplications.
    \item \textbf{Any $M > \max$ works}, not only $n$; choose $M = n$ only when all values are in $[0, n)$. For values up to $10^9$ no base fits in 64 bits twice --- use the cycle method.
    \item \textbf{One pass is impossible to shorten}: pass 2 cannot be merged into pass 1, because later iterations still need the old values of earlier cells.
\end{itemize}
\end{tricksbox}

\subsubsection*{5. Trace}

\begin{tracebox}[{Encoding on Sample 1, $n = 6$}]
\centering
\begin{tikzpicture}[x=1cm, y=1cm]
    \dpindex{ix}{2.75}{index}{0}{5}
    \dparray{e0}{2.1}{start}{0/acell, 2/acell, 1/acell, 5/acell, 3/acell, 4/acell}
    \dparray{e1}{0.7}{after $i=2$}{0/awrite, 8/awrite, 13/awrite, 5/acell, 3/acell, 4/acell}
    \dparray{e2}{-0.7}{after pass 1}{0/awrite, 8/awrite, 13/awrite, 29/awrite, 33/awrite, 22/awrite}
    \dparray{e3}{-2.1}{after pass 2}{0/akeep, 1/akeep, 2/akeep, 4/akeep, 5/akeep, 3/akeep}
    \node[note, anchor=west, text width=4.0cm, align=left] at (5.8,0.7) {blue = encoded: old $+\,6\cdot$new; e.g.\ $13 = 1 + 6 \cdot 2$};
    \node[note, anchor=west, text width=4.0cm, align=left] at (5.8,-0.7) {$29 = 5 + 6\cdot 4$, \ $33 = 3 + 6 \cdot 5$, \ $22 = 4 + 6 \cdot 3$};
    \node[note, anchor=west, text width=4.0cm, align=left] at (5.8,-2.1) {divide by 6: the answer $[0,1,2,4,5,3]$};
\end{tikzpicture}
\figcap{Snapshots of the array. At every moment each cell's remainder mod 6 is still its original value.}

\medskip
\small
\begin{tabular}{@{}cccccl@{}}
\toprule
$i$ & old $= \texttt{nums}[i] \% 6$ & cell read & new $= \texttt{nums}[\text{old}] \% 6$ & cell becomes & \textbf{array after pass-1 step} \\
\midrule
0 & 0 & 0 (holds 0) & 0 & $0 + 6\cdot0 = 0$ & $[0,2,1,5,3,4]$ \\
1 & 2 & 2 (holds 1) & 1 & $2 + 6\cdot1 = 8$ & $[0,8,1,5,3,4]$ \\
2 & 1 & 1 (holds 8) & $8 \% 6 = 2$ & $1 + 6\cdot2 = 13$ & $[0,8,13,5,3,4]$ \\
3 & 5 & 5 (holds 4) & 4 & $5 + 6\cdot4 = 29$ & $[0,8,13,29,3,4]$ \\
4 & 3 & 3 (holds 29) & $29 \% 6 = 5$ & $3 + 6\cdot5 = 33$ & $[0,8,13,29,33,4]$ \\
5 & 4 & 4 (holds 33) & $33 \% 6 = 3$ & $4 + 6\cdot3 = 22$ & $[0,8,13,29,33,22]$ \\
\bottomrule
\end{tabular}

\medskip
At $i = 2$, $4$ and $5$ the source cell was already encoded; \texttt{\% 6} recovers its old value --- exactly the step the naive in-place loop gets wrong.
\end{tracebox}

\begin{tracebox}[{The cycle method on Sample 1}]
\centering
\begin{tikzpicture}[x=1cm, y=1cm]
    \dpindex{cx}{3.05}{index}{0}{5}
    \dparray{k0}{2.4}{start}{0/acell, 2/acell, 1/acell, 5/acell, 3/acell, 4/acell}
    \dparray{k1}{1.2}{$s = 0$}{0/awrite, 2/acell, 1/acell, 5/acell, 3/acell, 4/acell}
    \dparray{k2}{0}{$s = 1$}{0/awrite, 1/awrite, 2/awrite, 5/acell, 3/acell, 4/acell}
    \dparray{k3}{-1.2}{$s = 3$}{0/akeep, 1/akeep, 2/akeep, 4/akeep, 5/akeep, 3/akeep}
    \node[note, anchor=west, text width=4.2cm, align=left] at (5.8,1.2) {fixed point: $0 \to 0$ keeps 0};
    \node[note, anchor=west, text width=4.2cm, align=left] at (5.8,0) {2-cycle $(1\,2)$: the values swap};
    \node[note, anchor=west, text width=4.2cm, align=left] at (5.8,-1.2) {3-cycle $(3\,5\,4)$: values shift; done};
\end{tikzpicture}
\figcap{The array after each cycle walk, with the marks already removed for readability (blue: rewritten in this row's walk).}

\medskip
\small
\begin{tabular}{@{}cll@{}}
\toprule
\textbf{start} $s$ & \textbf{cycle walked} & \textbf{array after the walk ($\sim v$ shown as $-v-1$)} \\
\midrule
0 & $0$ (fixed point) & $[-1,2,1,5,3,4]$ \\
1 & $1 \to 2$ & $[-1,-2,-3,5,3,4]$ \\
2, 3 & 2 already marked; $3 \to 5 \to 4$ & $[-1,-2,-3,-5,-6,-4]$ \\
4, 5 & already marked & unchanged \\
unmark & $v \leftarrow \sim v$ & $[0,1,2,4,5,3]$ \\
\bottomrule
\end{tabular}
\end{tracebox}

\subsubsection*{6. Test suite}

\begin{testbox}
\centering\small
\begin{tabular}{@{}lll@{}}
\toprule
\textbf{nums} & \textbf{Why it is interesting} & \textbf{Expected} \\
\midrule
$[0,2,1,5,3,4]$ & official Sample 1: three cycles & $[0,1,2,4,5,3]$ \\
$[5,0,1,2,3,4]$ & official Sample 2: one cycle of length 6 & $[4,5,0,1,2,3]$ \\
$[0]$ & $n = 1$ & $[0]$ \\
$[1,0]$ & a 2-cycle: applying it twice is the identity & $[0,1]$ \\
$[0,1,2,3]$ & identity & $[0,1,2,3]$ \\
$[1,2,0]$ & a 3-cycle: $p \circ p = p^{-1}$ & $[2,0,1]$ \\
$n = 1000$, reversed & largest encoded value $999 + 1000 \cdot 999 < 10^6$ & the identity \\
\bottomrule
\end{tabular}

\medskip
Both in-place programs reproduce the official samples and agree with the two-array version on 500 random permutations with $n \le 9$.
\end{testbox}

\subsubsection*{7. Common mistakes}

\begin{warningbox}[Mistake 1 --- in place without encoding]
\texttt{nums[i] = nums[nums[i]]} returns $[0,1,1,4,4,4]$ on Sample 1 (see the observations).
\end{warningbox}

\begin{warningbox}[Mistake 2 --- forgetting one of the two \texttt{\% n}]
\begin{lstlisting}[style=snippet]
nums[i] += n * nums[nums[i]];             // WRONG: the source may be encoded already
nums[i] += n * (nums[nums[i] % n] % n);   // RIGHT: old index, old value
\end{lstlisting}
The inner \texttt{\% n} is needed too, although cell $i$ itself is not yet encoded when it is read --- it costs nothing and keeps the formula symmetric and safe if the loop is ever reordered.
\end{warningbox}

\begin{warningbox}[Mistake 3 --- marking with the minus sign]
\texttt{nums[i] = -nums[i]} cannot mark the value 0, and a later test \texttt{nums[i] < 0} misses it. Use $\sim v$ or add an offset.
\end{warningbox}

\begin{warningbox}[Mistake 4 --- an encoding base that overflows]
The same trick on a problem with $n = 10^5$ and values below $n$ needs $n^2 \approx 10^{10}$: \texttt{int} overflows. Check $M^2$ against the type before using the encoding.
\end{warningbox}

\subsubsection*{8. Takeaways}

\begin{itemize}[leftmargin=*]
    \item When values are indices, a cell can be followed: $i \to a[i] \to a[a[i]]$; a permutation is a set of cycles.
    \item In-place rewriting fails when a new value needs old values of cells that were already rewritten; encoding $\text{old} + n\cdot\text{new}$ keeps both readable with \texttt{\%} and \texttt{/}.
    \item The encoding needs all values in $[0, M)$ and $M^2$ inside the type; the cycle walk needs one spare bit per cell.
    \item ``$\bigO{1}$ extra memory'' follow-ups are usually one of three tricks: reversal, cycles, or encoding.
\end{itemize}

\begin{variantsbox}
\variant{Concatenation of Array}{\exeasy\ \lc{Easy}\ \varlink{Concatenation of Array}{LeetCode 1929}{https://leetcode.com/problems/concatenation-of-array/}}
  {return the array of length $2n$ with $\texttt{ans}[i] = \texttt{ans}[i+n] = \texttt{nums}[i]$.}
  {model: index arithmetic $i \to i \bmod n$ for the output position $i \in [0, 2n)$.}
  {algorithm: \texttt{resize(2n)} and copy $a[i+n] = a[i]$; samples $\to [1,2,1,1,2,1]$, $[1,3,2,1,1,3,2,1]$.}
\variant{Shuffle the Array}{\exeasy\ \lc{Easy}\ \varlink{Shuffle the Array}{LeetCode 1470}{https://leetcode.com/problems/shuffle-the-array/}}
  {$[x_1..x_n, y_1..y_n] \to [x_1, y_1, x_2, y_2, \ldots]$; values $\le 10^3$.}
  {model: output position $i$ takes input position $i/2$ (even $i$) or $n + i/2$ (odd $i$).}
  {algorithm: in place with 10-bit packing: \texttt{a[i] |= (a[src] \& 1023) << 10}, then \texttt{>>= 10}; all three samples reproduced.}
\variant{Inverse permutation}{\varnone}
  {build $\texttt{inv}$ with $\texttt{inv}[\texttt{nums}[i]] = i$, in place.}
  {model: on every cycle the arrows are reversed; old $+\,n\cdot$new still works, but the write goes to cell $\texttt{nums}[i]$.}
  {algorithm: pass 1: \texttt{a[a[i] \% n] += n * i}; pass 2: divide by $n$. $\bigO{n}$ time, $\bigO{1}$ memory.}
\variant{Apply the permutation $k$ times}{\varnone}
  {$\texttt{ans}[i] = p^k(i)$ for a large $k$ (up to $10^{18}$).}
  {model: on a cycle of length $L$, $p^k$ moves every element $k \bmod L$ steps forward.}
  {algorithm: walk each cycle once, record it in a list, and read position $(t + k) \bmod L$; $\bigO{n}$ time.}
\end{variantsbox}
